\documentclass[11pt,a4paper]{article}
\usepackage[utf8]{inputenc}
\usepackage[T1]{fontenc}
\usepackage[french]{babel}
\usepackage{amsmath}
\usepackage{amssymb}
\usepackage{amsthm}
\usepackage{mathtools}
\usepackage{geometry}
\usepackage{enumitem}
\usepackage{xcolor}
\usepackage{titlesec}
\usepackage{fancyhdr}
\usepackage{booktabs}
\usepackage{array}
\usepackage{longtable}
\usepackage{tcolorbox}

% Configuration de la page
\geometry{margin=2.5cm}
\pagestyle{fancy}
\fancyhf{}
\fancyhead[L]{\leftmark}
\fancyhead[R]{Exercices de Révision - Correction}
\fancyfoot[C]{\thepage}

% Configuration des titres
\titleformat{\section}
{\Large\bfseries\color{blue!70!black}}
{}
{0em}
{}[\titlerule]

\titleformat{\subsection}
{\large\bfseries\color{blue!50!black}}
{}
{0em}
{}

\titleformat{\subsubsection}
{\normalsize\bfseries\color{blue!30!black}}
{}
{0em}
{}

% Boîtes colorées
\newtcolorbox{exemplebox}{
    colback=blue!5!white,
    colframe=blue!75!black,
    title=\textbf{Exemple},
    fonttitle=\bfseries
}

\newtcolorbox{definitionbox}{
    colback=green!5!white,
    colframe=green!75!black,
    title=\textbf{Définition},
    fonttitle=\bfseries
}

\newtcolorbox{astucebox}{
    colback=yellow!5!white,
    colframe=yellow!75!black,
    title=\textbf{Astuce},
    fonttitle=\bfseries
}

% Métadonnées
\title{Correction des Exercices de Révision\\
\large Analyse de Données - Test Chi-2 et ACP}
\author{EMSI Rabat - Filière 4IIR}
\date{2025/2026}

\begin{document}

\maketitle

\tableofcontents
\newpage

\section{Exercice 1 : Test d'Indépendance Chi-2}

\subsection{Énoncé}

On veut étudier la liaison entre les caractères : "être fumeur" (plus de 20 cigarettes par jour, pendant 10 ans) et "avoir un cancer de la gorge", sur une population de 1000 personnes, dont 500 sont atteintes d'un cancer de la gorge. 

Voici les résultats observés :

\begin{table}[h]
\centering
\begin{tabular}{|c|c|c|c|}
\hline
\textbf{Observé} & \textbf{Cancer} & \textbf{Non cancer} & \textbf{Marge} \\
\hline
Fumeur & 342 & 258 & 600 \\
\hline
Non fumeur & 158 & 242 & 400 \\
\hline
\textbf{Marge} & \textbf{500} & \textbf{500} & \textbf{1000} \\
\hline
\end{tabular}
\caption{Tableau des effectifs observés}
\end{table}

\textbf{Objectif :} Faire un test d'indépendance pour établir la liaison entre ces caractères.

\subsection{Méthodologie du Test Chi-2}

\begin{definitionbox}
Le test du Chi-2 (ou $\chi^2$) d'indépendance permet de tester si deux variables qualitatives sont indépendantes ou non.

\textbf{Étapes du test :}
\begin{enumerate}
    \item Formuler les hypothèses $H_0$ et $H_1$
    \item Choisir le risque d'erreur $\alpha$ (souvent 5\%)
    \item Calculer les effectifs attendus sous $H_0$
    \item Calculer la statistique de test $\chi_0^2$
    \item Déterminer le degré de liberté $\nu$
    \item Comparer avec la valeur critique $\chi_{\alpha,\nu}^2$
    \item Prendre une décision
\end{enumerate}
\end{definitionbox}

\subsection{Solution Détaillée}

\subsubsection{Étape 1 : Formulation des hypothèses}

\begin{itemize}
    \item \textbf{$H_0$} : Les deux variables sont \textbf{indépendantes}
    \item \textbf{$H_1$} : Les deux variables sont \textbf{dépendantes}
\end{itemize}

\subsubsection{Étape 2 : Choix du risque d'erreur}

Puisque $\alpha$ n'est pas fixé dans l'énoncé, on prend :
$$\alpha = 5\% = 0,05$$

\subsubsection{Étape 3 : Calcul des effectifs attendus sous $H_0$}

Sous l'hypothèse d'indépendance $H_0$, l'effectif attendu dans la case $(i,j)$ est calculé par :

$$E_{ij} = \frac{n_{i.} \times n_{.j}}{n}$$

où :
\begin{itemize}
    \item $n_{i.}$ : total de la ligne $i$
    \item $n_{.j}$ : total de la colonne $j$
    \item $n$ : effectif total
\end{itemize}

\textbf{Calculs détaillés :}

\begin{itemize}
    \item Fumeur + Cancer : $E_{11} = \frac{600 \times 500}{1000} = 300$
    \item Fumeur + Non cancer : $E_{12} = \frac{600 \times 500}{1000} = 300$
    \item Non fumeur + Cancer : $E_{21} = \frac{400 \times 500}{1000} = 200$
    \item Non fumeur + Non cancer : $E_{22} = \frac{400 \times 500}{1000} = 200$
\end{itemize}

\begin{table}[h]
\centering
\begin{tabular}{|c|c|c|c|}
\hline
\textbf{Attendu} & \textbf{Cancer} & \textbf{Non cancer} & \textbf{Marge} \\
\hline
Fumeur & 300 & 300 & 600 \\
\hline
Non fumeur & 200 & 200 & 400 \\
\hline
\textbf{Marge} & \textbf{500} & \textbf{500} & \textbf{1000} \\
\hline
\end{tabular}
\caption{Tableau des effectifs attendus sous $H_0$}
\end{table}

\subsubsection{Étape 4 : Calcul de la statistique de test}

La statistique du Chi-2 est calculée par :

$$\chi_0^2 = \sum_{i=1}^{r} \sum_{j=1}^{c} \frac{(O_{ij} - E_{ij})^2}{E_{ij}}$$

où :
\begin{itemize}
    \item $O_{ij}$ : effectif observé dans la case $(i,j)$
    \item $E_{ij}$ : effectif attendu dans la case $(i,j)$
    \item $r$ : nombre de lignes
    \item $c$ : nombre de colonnes
\end{itemize}

\textbf{Calcul détaillé :}

\begin{align*}
\chi_0^2 &= \frac{(342 - 300)^2}{300} + \frac{(258 - 300)^2}{300} + \frac{(158 - 200)^2}{200} + \frac{(242 - 200)^2}{200}\\
&= \frac{42^2}{300} + \frac{(-42)^2}{300} + \frac{(-42)^2}{200} + \frac{42^2}{200}\\
&= \frac{1764}{300} + \frac{1764}{300} + \frac{1764}{200} + \frac{1764}{200}\\
&= 5,88 + 5,88 + 8,82 + 8,82\\
&= 29,4
\end{align*}

$$\boxed{\chi_0^2 = 29,4}$$

\subsubsection{Étape 5 : Détermination du degré de liberté}

Le degré de liberté est calculé par :

$$\nu = (r - 1) \times (c - 1)$$

où :
\begin{itemize}
    \item $r = 2$ : nombre de modalités de la première variable (Fumeur/Non fumeur)
    \item $c = 2$ : nombre de modalités de la deuxième variable (Cancer/Non cancer)
\end{itemize}

$$\nu = (2 - 1) \times (2 - 1) = 1 \times 1 = 1$$

\subsubsection{Étape 6 : Valeur critique}

Dans la table du Chi-2, avec $\alpha = 0,05$ et $\nu = 1$, on trouve :

$$\chi_{0,05;1}^2 = 3,841$$

\subsubsection{Étape 7 : Décision}

\begin{astucebox}
\textbf{Règle de décision :}
\begin{itemize}
    \item Si $\chi_0^2 > \chi_{\alpha,\nu}^2$ : on \textbf{rejette} $H_0$ (les variables sont dépendantes)
    \item Si $\chi_0^2 \leq \chi_{\alpha,\nu}^2$ : on \textbf{accepte} $H_0$ (les variables sont indépendantes)
\end{itemize}
\end{astucebox}

\textbf{Application :}

$$\chi_0^2 = 29,4 > 3,841 = \chi_{0,05;1}^2$$

\subsection{Conclusion}

Puisque $\chi_0^2 = 29,4 > 3,841$, on \textbf{rejette} $H_0$ au seuil de signification de 5\%.

\textbf{Interprétation :} Il existe une \textbf{liaison significative} entre "être fumeur" et "avoir un cancer de la gorge". Les deux variables sont \textbf{dépendantes}.

\newpage

\section{Exercice 2 : Analyse en Composantes Principales (ACP)}

\subsection{Énoncé}

Nous allons analyser les données de ventes de deux produits dans plusieurs magasins. Le tableau ci-dessous présente les chiffres de ventes des produits A et B.

\begin{table}[h]
\centering
\begin{tabular}{|c|c|c|}
\hline
\textbf{Magasin} & \textbf{Ventes Produit A} & \textbf{Ventes Produit B} \\
\hline
1 & 2 & 2 \\
\hline
2 & 6 & 2 \\
\hline
3 & 6 & 4 \\
\hline
4 & 10 & 4 \\
\hline
\end{tabular}
\caption{Ventes des Produits A et B par Magasin}
\end{table}

\textbf{Questions :}
\begin{enumerate}
    \item Définir la matrice des données $X$.
    \item Calculer la moyenne et l'écart-type de chaque variable.
    \item Déterminer la matrice centrée et réduite $X_{cr}$.
    \item Déterminer la matrice des corrélations entre les variables $U$.
    \item Déterminer le polynôme caractéristique de la matrice de corrélation $P(\lambda)$.
    \item Calculer les valeurs propres.
    \item Déterminer les axes principaux et le pourcentage d'information expliqué par chaque axe.
    \item Déterminer les composantes principales.
    \item Déterminer la contribution de chaque individu à l'explication du phénomène, ainsi que celle des variables dans l'analyse.
    \item Calculer la contribution de chaque individu à la construction des axes principaux.
\end{enumerate}

\subsection{Méthodologie de l'ACP}

\begin{definitionbox}
L'Analyse en Composantes Principales (ACP) est une méthode d'analyse de données qui permet de :
\begin{itemize}
    \item Réduire la dimension des données
    \item Identifier les axes principaux de variation
    \item Visualiser les relations entre variables et individus
    \item Interpréter la structure des données
\end{itemize}

\textbf{Étapes principales :}
\begin{enumerate}
    \item Centrer et réduire les données
    \item Calculer la matrice de corrélation
    \item Déterminer les valeurs propres et vecteurs propres
    \item Calculer les composantes principales
    \item Analyser les contributions et qualités de représentation
\end{enumerate}
\end{definitionbox}

\subsection{Solution Détaillée}

\subsubsection{Question 1 : Matrice des données}

La matrice des données $X$ est définie par :

$$X = \begin{pmatrix}
2 & 2 \\
6 & 2 \\
6 & 4 \\
10 & 4
\end{pmatrix}$$

où :
\begin{itemize}
    \item Les lignes représentent les individus (magasins 1, 2, 3, 4)
    \item Les colonnes représentent les variables (Produit A, Produit B)
\end{itemize}

On note :
\begin{itemize}
    \item $X_1$ : la première colonne (Produit A)
    \item $X_2$ : la deuxième colonne (Produit B)
\end{itemize}

\subsubsection{Question 2 : Moyenne et écart-type}

\textbf{Moyenne de $X_1$ (Produit A) :}

$$\overline{X_1} = \frac{1}{n} \sum_{i=1}^{n} X_{1i} = \frac{1}{4}(2 + 6 + 6 + 10) = \frac{24}{4} = 6$$

\textbf{Moyenne de $X_2$ (Produit B) :}

$$\overline{X_2} = \frac{1}{n} \sum_{i=1}^{n} X_{2i} = \frac{1}{4}(2 + 2 + 4 + 4) = \frac{12}{4} = 3$$

\textbf{Variance de $X_1$ :}

$$s_1^2 = \frac{1}{n} \sum_{i=1}^{n} (X_{1i} - \overline{X_1})^2 = \frac{1}{4}[(2-6)^2 + (6-6)^2 + (6-6)^2 + (10-6)^2]$$

$$s_1^2 = \frac{1}{4}[16 + 0 + 0 + 16] = \frac{32}{4} = 8$$

\textbf{Écart-type de $X_1$ :}

$$s_1 = \sqrt{s_1^2} = \sqrt{8} = 2\sqrt{2}$$

\textbf{Variance de $X_2$ :}

$$s_2^2 = \frac{1}{n} \sum_{i=1}^{n} (X_{2i} - \overline{X_2})^2 = \frac{1}{4}[(2-3)^2 + (2-3)^2 + (4-3)^2 + (4-3)^2]$$

$$s_2^2 = \frac{1}{4}[1 + 1 + 1 + 1] = \frac{4}{4} = 1$$

\textbf{Écart-type de $X_2$ :}

$$s_2 = \sqrt{s_2^2} = \sqrt{1} = 1$$

\begin{astucebox}
\textbf{Rappel :} En ACP, on utilise la variance \textbf{empirique} (diviseur $n$) et non la variance corrigée (diviseur $n-1$).
\end{astucebox}

\subsubsection{Question 3 : Matrice centrée et réduite}

La matrice centrée et réduite $X_{cr}$ est obtenue en centrant (soustraire la moyenne) et en réduisant (diviser par l'écart-type) chaque variable.

L'élément $(i,j)$ de $X_{cr}$ est :

$$(X_{cr})_{ij} = \frac{X_{ij} - \overline{X_j}}{s_j}$$

\textbf{Calcul détaillé :}

$$X_{cr} = \begin{pmatrix}
\frac{2-6}{2\sqrt{2}} & \frac{2-3}{1} \\
\frac{6-6}{2\sqrt{2}} & \frac{2-3}{1} \\
\frac{6-6}{2\sqrt{2}} & \frac{4-3}{1} \\
\frac{10-6}{2\sqrt{2}} & \frac{4-3}{1}
\end{pmatrix} = \begin{pmatrix}
\frac{-4}{2\sqrt{2}} & -1 \\
0 & -1 \\
0 & 1 \\
\frac{4}{2\sqrt{2}} & 1
\end{pmatrix} = \begin{pmatrix}
-\sqrt{2} & -1 \\
0 & -1 \\
0 & 1 \\
\sqrt{2} & 1
\end{pmatrix}$$

$$\boxed{X_{cr} = \begin{pmatrix}
-\sqrt{2} & -1 \\
0 & -1 \\
0 & 1 \\
\sqrt{2} & 1
\end{pmatrix}}$$

\subsubsection{Question 4 : Matrice de corrélation}

La matrice de corrélation $U$ est calculée par :

$$U = \frac{1}{n} X_{cr}^T X_{cr}$$

\textbf{Calcul détaillé :}

$$X_{cr}^T = \begin{pmatrix}
-\sqrt{2} & 0 & 0 & \sqrt{2} \\
-1 & -1 & 1 & 1
\end{pmatrix}$$

$$U = \frac{1}{4} \begin{pmatrix}
-\sqrt{2} & 0 & 0 & \sqrt{2} \\
-1 & -1 & 1 & 1
\end{pmatrix} \begin{pmatrix}
-\sqrt{2} & -1 \\
0 & -1 \\
0 & 1 \\
\sqrt{2} & 1
\end{pmatrix}$$

$$U = \frac{1}{4} \begin{pmatrix}
(-\sqrt{2})(-\sqrt{2}) + 0 \cdot 0 + 0 \cdot 0 + \sqrt{2} \cdot \sqrt{2} & (-\sqrt{2})(-1) + 0 \cdot (-1) + 0 \cdot 1 + \sqrt{2} \cdot 1 \\
(-1)(-\sqrt{2}) + (-1) \cdot 0 + 1 \cdot 0 + 1 \cdot \sqrt{2} & (-1)(-1) + (-1)(-1) + 1 \cdot 1 + 1 \cdot 1
\end{pmatrix}$$

$$U = \frac{1}{4} \begin{pmatrix}
2 + 0 + 0 + 2 & \sqrt{2} + 0 + 0 + \sqrt{2} \\
\sqrt{2} + 0 + 0 + \sqrt{2} & 1 + 1 + 1 + 1
\end{pmatrix} = \frac{1}{4} \begin{pmatrix}
4 & 2\sqrt{2} \\
2\sqrt{2} & 4
\end{pmatrix}$$

$$U = \begin{pmatrix}
1 & \frac{\sqrt{2}}{2} \\
\frac{\sqrt{2}}{2} & 1
\end{pmatrix}$$

$$\boxed{U = \begin{pmatrix}
1 & \frac{\sqrt{2}}{2} \\
\frac{\sqrt{2}}{2} & 1
\end{pmatrix}}$$

\begin{astucebox}
La matrice de corrélation est toujours :
\begin{itemize}
    \item \textbf{Symétrique} : $U = U^T$
    \item Les éléments de la diagonale sont égaux à 1 (corrélation d'une variable avec elle-même)
    \item Les éléments hors diagonale sont les coefficients de corrélation entre variables
\end{itemize}
\end{astucebox}

\subsubsection{Question 5 : Polynôme caractéristique}

Le polynôme caractéristique $P(\lambda)$ est le déterminant de $(U - \lambda I_2)$ :

$$P(\lambda) = \det(U - \lambda I_2) = \det\begin{pmatrix}
1 - \lambda & \frac{\sqrt{2}}{2} \\
\frac{\sqrt{2}}{2} & 1 - \lambda
\end{pmatrix}$$

$$P(\lambda) = (1 - \lambda)^2 - \left(\frac{\sqrt{2}}{2}\right)^2 = (1 - \lambda)^2 - \frac{2}{4} = (1 - \lambda)^2 - \frac{1}{2}$$

$$P(\lambda) = 1 - 2\lambda + \lambda^2 - \frac{1}{2} = \lambda^2 - 2\lambda + \frac{1}{2}$$

On peut aussi factoriser :

$$P(\lambda) = (1 - \lambda)^2 - \left(\frac{\sqrt{2}}{2}\right)^2 = \left[(1 - \lambda) - \frac{\sqrt{2}}{2}\right]\left[(1 - \lambda) + \frac{\sqrt{2}}{2}\right]$$

$$P(\lambda) = \left(1 - \lambda - \frac{\sqrt{2}}{2}\right)\left(1 - \lambda + \frac{\sqrt{2}}{2}\right) = \left(\frac{2 - \sqrt{2}}{2} - \lambda\right)\left(\frac{2 + \sqrt{2}}{2} - \lambda\right)$$

$$\boxed{P(\lambda) = \left(\frac{2 - \sqrt{2}}{2} - \lambda\right)\left(\frac{2 + \sqrt{2}}{2} - \lambda\right)}$$

\subsubsection{Question 6 : Valeurs propres}

Les valeurs propres sont les racines du polynôme caractéristique :

$$P(\lambda) = 0 \Leftrightarrow \lambda = \frac{2 + \sqrt{2}}{2} \quad \text{ou} \quad \lambda = \frac{2 - \sqrt{2}}{2}$$

Dans l'ordre décroissant :

$$\lambda_1 = \frac{2 + \sqrt{2}}{2} \approx 1,707$$

$$\lambda_2 = \frac{2 - \sqrt{2}}{2} \approx 0,293$$

$$\boxed{\lambda_1 = \frac{2 + \sqrt{2}}{2}, \quad \lambda_2 = \frac{2 - \sqrt{2}}{2}}$$

\begin{astucebox}
\textbf{Propriétés des valeurs propres :}
\begin{itemize}
    \item La somme des valeurs propres = trace de $U$ = nombre de variables
    \item $\lambda_1 + \lambda_2 = 1 + 1 = 2$ (vérifié)
    \item Les valeurs propres sont toujours positives (matrice définie positive)
    \item $\lambda_1 \geq \lambda_2$ (ordre décroissant)
\end{itemize}
\end{astucebox}

\subsubsection{Question 7 : Axes principaux et pourcentage d'information}

\textbf{Calcul du premier axe principal $v_1$ :}

On cherche un vecteur propre $u_1$ associé à $\lambda_1$ tel que :

$$(U - \lambda_1 I_2) u_1 = 0$$

Avec $\lambda_1 = \frac{2 + \sqrt{2}}{2}$, on a :

$$U - \lambda_1 I_2 = \begin{pmatrix}
1 - \frac{2 + \sqrt{2}}{2} & \frac{\sqrt{2}}{2} \\
\frac{\sqrt{2}}{2} & 1 - \frac{2 + \sqrt{2}}{2}
\end{pmatrix} = \begin{pmatrix}
\frac{2 - 2 - \sqrt{2}}{2} & \frac{\sqrt{2}}{2} \\
\frac{\sqrt{2}}{2} & \frac{2 - 2 - \sqrt{2}}{2}
\end{pmatrix} = \begin{pmatrix}
-\frac{\sqrt{2}}{2} & \frac{\sqrt{2}}{2} \\
\frac{\sqrt{2}}{2} & -\frac{\sqrt{2}}{2}
\end{pmatrix}$$

Soit $u_1 = \begin{pmatrix} x \\ y \end{pmatrix}$, on résout :

$$\begin{pmatrix}
-\frac{\sqrt{2}}{2} & \frac{\sqrt{2}}{2} \\
\frac{\sqrt{2}}{2} & -\frac{\sqrt{2}}{2}
\end{pmatrix} \begin{pmatrix} x \\ y \end{pmatrix} = \begin{pmatrix} 0 \\ 0 \end{pmatrix}$$

$$\Leftrightarrow \begin{cases}
-\frac{\sqrt{2}}{2}x + \frac{\sqrt{2}}{2}y = 0 \\
\frac{\sqrt{2}}{2}x - \frac{\sqrt{2}}{2}y = 0
\end{cases} \Leftrightarrow x = y$$

Donc $u_1 = \begin{pmatrix} x \\ x \end{pmatrix} = x \begin{pmatrix} 1 \\ 1 \end{pmatrix}$

En prenant $x = 1$, on a $u_1 = \begin{pmatrix} 1 \\ 1 \end{pmatrix}$

La norme de $u_1$ est : $\|u_1\| = \sqrt{1^2 + 1^2} = \sqrt{2}$

Le vecteur normalisé est :

$$v_1 = \frac{1}{\|u_1\|} u_1 = \frac{1}{\sqrt{2}} \begin{pmatrix} 1 \\ 1 \end{pmatrix}$$

$$\boxed{v_1 = \frac{1}{\sqrt{2}} \begin{pmatrix} 1 \\ 1 \end{pmatrix}}$$

\textbf{Calcul du deuxième axe principal $v_2$ :}

Avec $\lambda_2 = \frac{2 - \sqrt{2}}{2}$, on a :

$$U - \lambda_2 I_2 = \begin{pmatrix}
\frac{\sqrt{2}}{2} & \frac{\sqrt{2}}{2} \\
\frac{\sqrt{2}}{2} & \frac{\sqrt{2}}{2}
\end{pmatrix}$$

Soit $u_2 = \begin{pmatrix} x \\ y \end{pmatrix}$, on résout :

$$\begin{pmatrix}
\frac{\sqrt{2}}{2} & \frac{\sqrt{2}}{2} \\
\frac{\sqrt{2}}{2} & \frac{\sqrt{2}}{2}
\end{pmatrix} \begin{pmatrix} x \\ y \end{pmatrix} = \begin{pmatrix} 0 \\ 0 \end{pmatrix}$$

$$\Leftrightarrow \frac{\sqrt{2}}{2}(x + y) = 0 \Leftrightarrow x = -y$$

Donc $u_2 = \begin{pmatrix} x \\ -x \end{pmatrix} = x \begin{pmatrix} 1 \\ -1 \end{pmatrix}$

En prenant $x = 1$, on a $u_2 = \begin{pmatrix} 1 \\ -1 \end{pmatrix}$

La norme de $u_2$ est : $\|u_2\| = \sqrt{1^2 + (-1)^2} = \sqrt{2}$

Le vecteur normalisé est :

$$v_2 = \frac{1}{\|u_2\|} u_2 = \frac{1}{\sqrt{2}} \begin{pmatrix} 1 \\ -1 \end{pmatrix}$$

$$\boxed{v_2 = \frac{1}{\sqrt{2}} \begin{pmatrix} 1 \\ -1 \end{pmatrix}}$$

\textbf{Pourcentage d'information expliqué :}

Le pourcentage d'information expliqué par l'axe $k$ est :

$$I_k = \frac{\lambda_k}{\sum_{i=1}^{p} \lambda_i} \times 100\%$$

où $p$ est le nombre de variables (ici $p = 2$).

\textbf{Premier axe :}

$$I_1 = \frac{\lambda_1}{\lambda_1 + \lambda_2} = \frac{\frac{2 + \sqrt{2}}{2}}{2} = \frac{2 + \sqrt{2}}{4} = \frac{1}{2} + \frac{\sqrt{2}}{4} \approx 0,854 = 85,4\%$$

\textbf{Deuxième axe :}

$$I_2 = \frac{\lambda_2}{\lambda_1 + \lambda_2} = \frac{\frac{2 - \sqrt{2}}{2}}{2} = \frac{2 - \sqrt{2}}{4} = \frac{1}{2} - \frac{\sqrt{2}}{4} \approx 0,146 = 14,6\%$$

$$\boxed{I_1 \approx 85,4\%, \quad I_2 \approx 14,6\%}$$

\begin{astucebox}
Le premier axe principal explique la majeure partie de la variance (85,4\%), ce qui indique une forte corrélation entre les deux variables.
\end{astucebox}

\subsubsection{Question 8 : Composantes principales}

Les composantes principales sont calculées par :

$$c^k = X_{cr} v_k$$

où $c^k$ est la $k$-ème composante principale.

\textbf{Première composante principale $c^1$ :}

$$c^1 = X_{cr} v_1 = \begin{pmatrix}
-\sqrt{2} & -1 \\
0 & -1 \\
0 & 1 \\
\sqrt{2} & 1
\end{pmatrix} \frac{1}{\sqrt{2}} \begin{pmatrix} 1 \\ 1 \end{pmatrix} = \frac{1}{\sqrt{2}} \begin{pmatrix}
-\sqrt{2} - 1 \\
0 - 1 \\
0 + 1 \\
\sqrt{2} + 1
\end{pmatrix}$$

$$c^1 = \frac{1}{\sqrt{2}} \begin{pmatrix}
-(\sqrt{2} + 1) \\
-1 \\
1 \\
\sqrt{2} + 1
\end{pmatrix} = \begin{pmatrix}
-\frac{\sqrt{2} + 1}{\sqrt{2}} \\
-\frac{1}{\sqrt{2}} \\
\frac{1}{\sqrt{2}} \\
\frac{\sqrt{2} + 1}{\sqrt{2}}
\end{pmatrix}$$

$$\boxed{c^1 = \begin{pmatrix}
-\frac{\sqrt{2} + 1}{\sqrt{2}} \\
-\frac{1}{\sqrt{2}} \\
\frac{1}{\sqrt{2}} \\
\frac{\sqrt{2} + 1}{\sqrt{2}}
\end{pmatrix}}$$

\textbf{Deuxième composante principale $c^2$ :}

$$c^2 = X_{cr} v_2 = \begin{pmatrix}
-\sqrt{2} & -1 \\
0 & -1 \\
0 & 1 \\
\sqrt{2} & 1
\end{pmatrix} \frac{1}{\sqrt{2}} \begin{pmatrix} 1 \\ -1 \end{pmatrix} = \frac{1}{\sqrt{2}} \begin{pmatrix}
-\sqrt{2} + 1 \\
0 + 1 \\
0 - 1 \\
\sqrt{2} - 1
\end{pmatrix}$$

$$c^2 = \frac{1}{\sqrt{2}} \begin{pmatrix}
1 - \sqrt{2} \\
1 \\
-1 \\
\sqrt{2} - 1
\end{pmatrix} = \begin{pmatrix}
\frac{1 - \sqrt{2}}{\sqrt{2}} \\
\frac{1}{\sqrt{2}} \\
-\frac{1}{\sqrt{2}} \\
\frac{\sqrt{2} - 1}{\sqrt{2}}
\end{pmatrix}$$

$$\boxed{c^2 = \begin{pmatrix}
\frac{1 - \sqrt{2}}{\sqrt{2}} \\
\frac{1}{\sqrt{2}} \\
-\frac{1}{\sqrt{2}} \\
\frac{\sqrt{2} - 1}{\sqrt{2}}
\end{pmatrix}}$$

\subsubsection{Question 9 : Contributions}

\textbf{Contribution des individus à l'inertie totale :}

La contribution de l'individu $i$ à l'axe $k$ est :

$$CTR_i^k = \frac{(c_i^k)^2}{n \lambda_k}$$

où $c_i^k$ est la coordonnée de l'individu $i$ sur l'axe $k$.

\textbf{Contribution des variables à l'axe $k$ :}

La contribution de la variable $j$ à l'axe $k$ est donnée par le carré de la coordonnée de la variable sur l'axe, c'est-à-dire le carré de l'élément $(j,k)$ du vecteur propre $v_k$.

Pour $v_1 = \frac{1}{\sqrt{2}} \begin{pmatrix} 1 \\ 1 \end{pmatrix}$ :
\begin{itemize}
    \item Variable 1 : $\left(\frac{1}{\sqrt{2}}\right)^2 = \frac{1}{2} = 50\%$
    \item Variable 2 : $\left(\frac{1}{\sqrt{2}}\right)^2 = \frac{1}{2} = 50\%$
\end{itemize}

Pour $v_2 = \frac{1}{\sqrt{2}} \begin{pmatrix} 1 \\ -1 \end{pmatrix}$ :
\begin{itemize}
    \item Variable 1 : $\left(\frac{1}{\sqrt{2}}\right)^2 = \frac{1}{2} = 50\%$
    \item Variable 2 : $\left(-\frac{1}{\sqrt{2}}\right)^2 = \frac{1}{2} = 50\%$
\end{itemize}

\begin{astucebox}
Les deux variables contribuent également à chaque axe (50\% chacune), ce qui est cohérent avec le fait que les deux variables ont la même variance après réduction.
\end{astucebox}

\subsubsection{Question 10 : Contribution des individus aux axes}

\textbf{Contribution au premier axe ($\lambda_1 = \frac{2 + \sqrt{2}}{2}$) :}

Pour chaque individu $i$ :

$$CTR_i^1 = \frac{(c_i^1)^2}{n \lambda_1} = \frac{(c_i^1)^2}{4 \times \frac{2 + \sqrt{2}}{2}} = \frac{(c_i^1)^2}{2(2 + \sqrt{2})}$$

\textbf{Calculs détaillés :}

\begin{itemize}
    \item Individu 1 : $CTR_1^1 = \frac{\left(-\frac{\sqrt{2} + 1}{\sqrt{2}}\right)^2}{2(2 + \sqrt{2})} = \frac{(\sqrt{2} + 1)^2}{2 \times 2(2 + \sqrt{2})} = \frac{(3 + 2\sqrt{2})}{4(2 + \sqrt{2})}$
    \item Individu 2 : $CTR_2^1 = \frac{\left(-\frac{1}{\sqrt{2}}\right)^2}{2(2 + \sqrt{2})} = \frac{1}{4(2 + \sqrt{2})}$
    \item Individu 3 : $CTR_3^1 = \frac{\left(\frac{1}{\sqrt{2}}\right)^2}{2(2 + \sqrt{2})} = \frac{1}{4(2 + \sqrt{2})}$
    \item Individu 4 : $CTR_4^1 = \frac{\left(\frac{\sqrt{2} + 1}{\sqrt{2}}\right)^2}{2(2 + \sqrt{2})} = \frac{(3 + 2\sqrt{2})}{4(2 + \sqrt{2})}$
\end{itemize}

\textbf{Contribution au deuxième axe ($\lambda_2 = \frac{2 - \sqrt{2}}{2}$) :}

$$CTR_i^2 = \frac{(c_i^2)^2}{n \lambda_2} = \frac{(c_i^2)^2}{4 \times \frac{2 - \sqrt{2}}{2}} = \frac{(c_i^2)^2}{2(2 - \sqrt{2})}$$

\begin{astucebox}
\textbf{Note importante :} Les contributions des individus aux axes permettent d'identifier quels individus sont les plus représentatifs de chaque axe principal. Les individus avec les contributions les plus élevées sont ceux qui "tirent" le plus l'axe dans leur direction.
\end{astucebox}

\newpage

\section{Résumé et Formules Clés}

\subsection{Test Chi-2 d'Indépendance}

\begin{itemize}
    \item \textbf{Hypothèses :} $H_0$ : indépendance, $H_1$ : dépendance
    \item \textbf{Effectifs attendus :} $E_{ij} = \frac{n_{i.} \times n_{.j}}{n}$
    \item \textbf{Statistique de test :} $\chi_0^2 = \sum_{i,j} \frac{(O_{ij} - E_{ij})^2}{E_{ij}}$
    \item \textbf{Degré de liberté :} $\nu = (r-1)(c-1)$
    \item \textbf{Décision :} Rejeter $H_0$ si $\chi_0^2 > \chi_{\alpha,\nu}^2$
\end{itemize}

\subsection{Analyse en Composantes Principales (ACP)}

\begin{itemize}
    \item \textbf{Matrice centrée-réduite :} $(X_{cr})_{ij} = \frac{X_{ij} - \overline{X_j}}{s_j}$
    \item \textbf{Matrice de corrélation :} $U = \frac{1}{n} X_{cr}^T X_{cr}$
    \item \textbf{Valeurs propres :} Solutions de $\det(U - \lambda I) = 0$
    \item \textbf{Axes principaux :} Vecteurs propres normalisés de $U$
    \item \textbf{Composantes principales :} $c^k = X_{cr} v_k$
    \item \textbf{Pourcentage d'information :} $I_k = \frac{\lambda_k}{\sum \lambda_i} \times 100\%$
    \item \textbf{Contribution individu :} $CTR_i^k = \frac{(c_i^k)^2}{n \lambda_k}$
\end{itemize}

\vspace{2cm}

\begin{center}
\textbf{Fin du document}
\end{center}

\end{document}

